\section{Emergence criteria: causality, isotropy, Nielsen--Ninomiya}
\label{p:criteria}

A consistent ladder is not enough on its own: one must also state what
counts as the relativistic limit being \emph{well posed}. We collect
three operational criteria, each verified by an explicit test.

\subsection{Effective light cone and causality}
Given a Dirac-like dispersion $E(\vp)^{2} = v^{2}\lvert\vp\rvert^{2}+m^{2}$, the
group velocity satisfies
\begin{equation}
v_{g}^{2}(\vp) - v^{2}
   \;=\; -\,\frac{m^{2}\,v^{2}}{v^{2}\lvert\vp\rvert^{2}+m^{2}} \;\le\; 0 ,
\label{p:eq:vg-bound}
\end{equation}
with saturation in the massless limit. This is the algebraic basis of
the effective light cone. The associated effective metric
$g^{\mu\nu}_{\mathrm{eff}} = \mathrm{diag}(-1,1/v^{2},1/v^{2},1/v^{2})$
has Lorentzian signature $(-,+,+,+)$ and determinant $-1/v^{6}$, both
verified symbolically (\texttt{validators/causality.py}). The
discrete-time analogue obeys $\cos(\varepsilon\,\Delta t) =
1-(v^{2}\vk^{2}+m^{2})\,\Delta t^{2}/2 + O(\Delta t^{3})$, recovering the
continuous dispersion at small step
(\texttt{validators/causalidad\_continuo\_vs\_discreto.py}).

\subsection{Isotropy as effective Lorentz}
For an anisotropic effective Hamiltonian
$\Hop_{\mathrm{eff}} = \sum_{a} v_{a}\,\sigma_{a}\,p_{a}$,
the squared dispersion
$\Hop_{\mathrm{eff}}^{2}=\sum_{a}v_{a}^{2}\,p_{a}^{2}\,\bbI$ reads as an
effective Riemannian metric
$g^{ab}_{\mathrm{eff}} = \mathrm{diag}(v_{x}^{2},v_{y}^{2},v_{z}^{2})$;
isotropy (a single effective light cone) is the condition
$v_{x}=v_{y}=v_{z}$. This is encoded as a SMT problem and verified by
\texttt{validators/isotropy.py}, which both confirms the existence of
isotropic solutions and the unsatisfiability of any combined
isotropy-with-anisotropy claim. For position-dependent velocities
$v_{a}(x)$ (variable tetrad) the same calculation goes through pointwise
(\texttt{validators/triada.py}).

\subsection{Chirality balance: Nielsen--Ninomiya without torus}
On any compact periodic lattice, the total chirality summed over band
nodes vanishes. Translated to combinatorial SMT, the statement is: for
$n$ even, there exists an assignment of charges $\chi_{i}\in\{\pm 1\}$
with $\sum_{i}\chi_{i}=0$; for $n$ odd, no such assignment exists.
Moreover, no uniform-chirality configuration is allowed when the sum is
zero. This is precisely what
\texttt{validators/nielsen\_ninomiya.py} verifies.

A stronger statement holds without invoking periodicity. For a finite
bipartite graph $G=(A\sqcup B,E)$, the chirality matrix
$\Gamma=\mathrm{diag}(+\bbI_{A},-\bbI_{B})$ anticommutes with the
adjacency, and the nullity is lower-bounded by $\lvert\lvert A\rvert-\lvert B\rvert\rvert$. On
a maximum matching the bound saturates. These facts, contained in
\texttt{validators/rama\_estructural.py}, are exactly the discrete index
theorem in the absence of reciprocal-space periodicity --- a key
ingredient of the structural branch of Sec.~\ref{p:open-problems}.

\subsection{Wilson term and the selection of the light Dirac}
On the 3D Wilson--Dirac lattice of Eq.~\eqref{p:eq:Wilson3D}, the
seven Brillouin-zone corners other than $\vk=0$ acquire Wilson masses
$m+2r, m+4r, m+6r$ depending on their distance from the origin in $\pi/a$
units, while $\vk=0$ retains mass $m$. This is the lattice mechanism
that lifts the naive fermion doublers and isolates a single light Dirac
in the infrared (\texttt{validators/wilson\_subsector.py},
\texttt{validators/protoespacio\_global.py}).
